%% maestro2tex -- generated table; do not edit by hand.
\begin{table}[htbp]
  \centering
  \caption[{Amplifier stability in unity-gain feedback by process corner.}]{Amplifier stability in unity-gain feedback by process corner. Unity-feedback bench (Figure~\ref{fig:tb-unityfb}) at $v_{\mathrm{cm}} = \SI{1.25}{\volt}$, $C_L = \SI{5}{\pico\farad}$, $C_c = \SI{500}{\femto\farad}$, \SI{2.5}{\volt} supply; the typical point is at \SI{37}{\celsius} and the four skew corners at \SI{25}{\celsius}. Measured through a probe in the feedback wire, so the quantity is the loop gain of the unity-gain configuration and not the amplifier loaded by an application network. Gain margin is read at the first crossing of zero loop phase above the unity-gain frequency.}
  \label{tab:ampab-unity}
  \begin{tabular}{lrrrrrrr}
    \toprule
    Parameter & Unit & tt & ff & fs & sf & ss & Spread \\
    \midrule
    Open-loop gain $A_{\mathrm{OL}}$ & \si{\decibel} & 82.74 & 80.28 & 83.30 & 83.31 & 86.75 & 6.47 \\
    Unity-gain frequency $f_{\mathrm{u}}$ & \si{\mega\hertz} & 22.474 & 22.942 & 22.621 & 22.587 & 22.358 & 0.584 \\
    Phase margin & \si{\degree} & 73.43 & 74.26 & 73.81 & 73.43 & 72.76 & 1.49 \\
    Gain margin & \si{\decibel} & 8.81 & 9.75 & 8.73 & 8.95 & 8.12 & 1.63 \\
    \bottomrule
  \end{tabular}
\end{table}
